% Q2V3: Series RLC Step Response, Critically Damped — Exam diagram
% Vs connected at t=0 via closing switch, series R-L-C, measure v_C(t)
\documentclass[border=10pt]{standalone}
\usepackage[american voltages, american currents]{circuitikz}
\usepackage{amsmath}
\renewcommand{\familydefault}{\sfdefault}

\begin{document}
\begin{circuitikz}[line width=0.8pt, every node/.style={font=\sffamily}]

% === Voltage source (left, vertical) ===
\draw (0,5) to[V, v=$V_s$] (0,0);

% === Closing switch on top wire ===
\draw (0,5) to[closing switch] (3,5);

% Switch labels (separate nodes, spaced vertically)
\node[above, font=\sffamily\itshape\small] at (1.5,5.6) {SW1};
\node[above, font=\sffamily\footnotesize] at (1.5,5.2) {$t{=}0$ (closes)};

% === Series R-L from switch to right side ===
\draw (3,5) to[R, l=$R$, i=$i(t)$] (6,5)
            -- (6,5) to[L, l=$L$] (6,2.5)
            to[C, l=$C$, v^=$v_C(t)$] (6,0);

% === Bottom return wire ===
\draw (6,0) -- (0,0);

% === Junction dots ===
\fill (0,0) circle (2pt);
\fill (0,5) circle (2pt);
\fill (6,5) circle (2pt);

% === Ground reference (optional clarity) ===
\node[ground] at (3,0) {};
\draw (0,0) -- (3,0);

\end{circuitikz}
\end{document}
